\section{实验原理}

\subsection{差分路线}
$i$轮差分路线是指对输入对$(X,X')$，在$i$轮加密过程中满足对应第$j$轮的一对输出值$(Y,Y')(1
le j \le i)$，满足$\delta X = X \oplus X' = \beta_0$，$\delta Y =  \oplus Y' = beta_j(1 \le j \le i)$，则称
$$
 \beta_0  \to \beta_1 \to \dots \to \beta_i 
$$
为一条$i$轮差分路线。假设密码算法的输入$X$和各轮密钥$K_i$的取值是相互独立且均匀分布的，对于满足输入差分为$beta_0$、输出差分为$\beta_i$(显然满足输入差分为$\beta_0$的明文不只一对)的$i$轮差分路线概率为$DP(\beta_0 \to \beta_i)$,其概率为各轮差分路线概率的乘积
$$
\text{DP}(\beta_0 \to \beta_1 \to \dots \to \beta_i)  = \prod_{(j=1)}^{i}\text{DP}(\beta_{j-1} \to \beta_j)
$$
注意这里差分路线是精确记录了包括起始、最终结果和中间每一种状态的一条路线。影响上述$i$轮差分路线的概率的主要因素是输入差分非零的$S$盒的个数以及输入差分非零的$S$盒对应的输出差分(因为当前轮$S$盒的输出即为下一轮$S$盒的输入)，其中输入非零的$S$盒被称为活跃$S$盒。

\subsection{差分}
$i$轮差分是指给定了输入差分$\beta_0$和$i$轮加密后的输出差分$\beta_i$，满足上述条件的称为$i$轮差分。与$i$轮差分路线相比，它不指定加密的中间状态。因此，对于$i$轮差分的概率计算如下：假设满足输入差分为$\beta_0$且$i$轮之后
输出差分为$\beta_i$的$i$轮差分路线共$s$条，则:
$$
\begin{aligned}
    \text{DP}(\beta_0 \xrightarrow{\text{i轮}} \beta_i) &= \sum_{t=1}^s \text{DP}( \beta_0 \xrightarrow{\text{1轮}} \beta_1^t \to \dots \beta_{i-1}^t \xrightarrow{\text{1轮}} \beta_i)\\
                                                        &= \sum_{t=1}^s (\prod_{j=1}^i \text{DP} (\beta_{j-1}^t \xrightarrow{\text{1轮}} \beta_j^t))
\end{aligned}
$$
也就是说$i$轮差分的概率是由所有满足输入和输出条件的$i$轮差分路线共同组成的，它的概率是要大于等于单条差分路线的概率的，而且想要精确计算差分的概率，还是要回到计算每一条差分路线的概率上去。

\subsection{CipherFour算法}

\subsection{利用线性规划寻找差分路线}
定义\(B_r(\Delta)\)为经过 \(r\) 轮后，到达差分状态 \(\Delta\) 的所有路线中，概率最大的一条的概率。
对于输入差分，把输入差分为0直接丢掉，因为输入差分为0之后的概率都是1，只能找到全0的路线，定义
\[
B_0(\Delta)=\begin{cases}
1,&\Delta\neq0,\\
0,&\Delta=0.
\end{cases}
\]

沿用书上的符号，假设每一轮的差分为\(\beta_{r}\)

假设已经知道了\(B_{r-1}(\beta_{r-1}),\),现在要计算\(B_r(\beta_{r})\)

从前驱 \(\beta_{r-1}\) 到达 \(\beta_{r}\) 的最佳完整路线概率为
\[
\Pr(\beta_{r-1}\rightarrow \beta_{r})=B_{r-1}(\beta_{r-1})\cdot \Pr(\beta_{r-1},\beta_{r}).
\]

然后再遍历所有前驱，就有

\[
B_r(\beta_{r}) = \max_{\beta_{r-1}}\{B_{r-1}(\beta_{r-1})\cdot \Pr(\beta_{r-1},\beta_{r})\}
\]

这在代码上看是比较简单的，只要遍历所有\(\beta_{r-1}\)就行，\(B_{r-1}(\beta_{r-1}),Pr(\beta_{r-1},\beta_{r})\)知道\(\beta_{r-1}\)后可以直接查表

但是理论上的东西，并不总是可行的，要知道每个\(\beta_{r-1}\)可是有$2^{16}$个，看上去也不是很多，但是注意这里的时间复杂度是在指数上的，也就是\(O(2^{16r})\)

那么如果是4轮，就会达到$2^{64}$,那这就不是一个小数目了，所以要拆分成4个S盒